\section{Discussion and Conclusion}
\textbf{Limitations and Future Work}. A limitation of our study is the simplification of texture generation to the main color style prediction due to the complexities of detailed texture synthesis. Extending this work to generate complete 3D textures is a key focus for future research. Moreover, given the substantial individual variations in EEG, future work should extend to enhance cross-subject generalization.\\ 
\textbf{Conclusion}. We explore a new task of reconstructing colored 3D objects from EEG signals, which is challenging but holds considerable importance for understanding the brain's mechanisms for real-time 3D perception. To facilitate this task, we develop the EEG-3D dataset, which integrates multimodal data and extensive EEG recordings. This dataset addresses the scarcity of EEG-3D object pairings, providing a valuable resource for future research in this domain. Furthermore, we propose a new framework, Neuro-3D, for extracting EEG-based visual features and reconstructing 3D objects. Neuro-3D leverages a diffusion-based 3D decoder for shape and color generation, conditioned on adaptively fused EEG features captured under static and dynamic 3D stimuli.
Extensive experiments demonstrate the feasibility of decoding 3D information from EEG signals and confirms the alignment between EEG visual decoding and biological visual perception mechanisms.